% SPDX-License-Identifier: CC-BY-SA-4.0
% Author: Matthieu Perrin
% Part: 
% Section: 
% Sub-section: 
% Frame: 

\begingroup

\begin{frame}{Caractérisation}

  \begin{block}{Point de linéarisation}
    Un objet partagé est linéarisable si, et seulement si : \\
    le résultat observable des opérations sur l'objet partagé d'une exécution est le même que si
    \alert{chaque opération avait lieu en un point atomique
      situé entre son début et sa fin.}
  \end{block}

  \begin{exampleblock}{Exemple}
    \begin{center}
      \begin{tikzpicture}[y=7mm]
        \draw[process, example] (0,4) node[left]{$counter$} to (8,4);
        \draw[process]          (0,3) node[left]{$T_1$}     to (8,3);
        \draw[process]          (0,2) node[left]{$T_2$}     to (8,2);

        \draw[example, lin point] (3,2) -- (3,4) node{$\bullet$};
        \draw[example, lin point] (6,3) -- (6,4) node{$\bullet$};

        \node[example, operation] at (3,2)   {$increment()$ };
        \node[example, operation] at (6,3) {$get()$};
      \end{tikzpicture}
    \end{center}
  \end{exampleblock}

  \begin{alertblock}{Propriétés}
    \begin{itemize}
    \item La linéarisabilité est composable
    \item La linéarisabilité implique la cohérence séquentielle
    \item Si chaque objet est linéarisable, l'exécution est SC
    \end{itemize}
  \end{alertblock}

\end{frame}

\endgroup
\endinput
